\scanpage{17}
% archive foliation: 16
$X$ courbe analytique connexe compacte genre $g\ge\unsure{2|1}$

$M_g^{\unsure{\xi|\zeta}}(X)$ espace modulaire pour les fibrés projectifs analytiques à connexion plate sur $X$, ``sans automorphismes''. C'est l'espace analytique associé à un schéma (de t.f. ?) sur $\mathbb C$, \deleted{mais} lisse de dim $6g-6$, mais on le regarde pt de vue analytique. Il est alors can. isomorphe (une fois choisi pt base $x\in X$) à l'espace des représentations
\diagram{\begin{tikzcd}
\pi\arrow[r,"\rho"]\arrow[d,phantom,"\Vert" description]&GP(1,\mathbb C)=H\\
\pi_1(X,x)
\end{tikzcd}}
\deleted{\unsure{telles}} que $\operatorname{Cent}(\rho)=\{e\}$, \deleted{\unsure{\#}} \unsure{quotienté} par $H$ opérant par composition avec aut. intérieurs
\[
M_g^{\unsure{\xi|\zeta}}(X)\simeq {}_H\!\backslash\uline{\operatorname{Hom}}^!_{\operatorname{gr.}\deleted{\unsure{\mathrm{cont}}}}(\pi,H)\deleted{\unsure{/H}}
\]
où le $!$ désigne le sous-ens. des $\rho\in\uline{\operatorname{Hom}}_{\operatorname{gr}}(\pi,H)$ dont le stabilisateur dans $H$ est réduit à $1$.

[\uline{NB} Il est immédiat de faire de $\uline{\operatorname{Hom}}_{\operatorname{gr}}(\pi,H)$ un espace algébrique (de façon que pour $g\in\pi$, les morphismes $\rho\longmapsto\rho(g)$ soient algébriques), de façon que $H$ y opère algébriquement. Mais l'iso précédent \uline{n'est pas} \inserted{\uline{sans doute}} algébrique !]
\margin{$?\ \big|$}
L'interprétation ``géom.'' de $\operatorname{Hom}^!$ lui-même est : on classifie les couples $(P,c,\varphi)$, $(P,c)$ un
